\section{Complete data and exact certificates}
\label{appA}\label{app:data}
This appendix gives the complete input and the exact certificates for the counterexample of the
preceding section; all arithmetic is over the rationals, with no rounding.
\subsection{Input and derived quantities}
\label{appA:input}\label{appA:derived}
Take $n=8$, $p=2$, $q=4$ rows and $L_0$ empty ($r_0=0$), with
\begin{equation}
\label{appA:eq:C}
C=\begin{bmatrix}1&0&0&0&0&0&0&0\\0&1&0&0&0&0&0&0\end{bmatrix},
\end{equation}
\begin{equation}
\label{appA:eq:A}
A=\begin{bmatrix}
0&0&1&0&0&0&0&0\\0&0&0&1&0&0&0&0\\
1&-2&1&0&-1&1&1&-1\\-1&2&1&0&2&-2&-1&2\\
0&0&0&0&1&0&0&0\\0&0&0&0&0&1&0&0\\
0&0&0&0&0&0&2&0\\0&0&0&0&0&0&0&3
\end{bmatrix}.
\end{equation}
For the data of \eqref{appA:eq:C},
\begin{equation}\label{appA:eq:rank}
\rank\begin{bmatrix}C\\ CA\end{bmatrix}=4,\qquad k=4,\qquad \rho=4,\qquad \omega=3,
\end{equation}
and a basis of $K_0$, the compression, the characteristic polynomial and the complementary block are
\begin{gather}
\label{appA:eq:N}
N=\begin{bmatrix}
0&0&0&0\\0&0&0&0\\0&0&0&0\\0&0&0&0\\
1&0&0&0\\0&1&0&0\\0&0&1&0\\0&0&0&1
\end{bmatrix},\\
M=(N^{\mathsf T}N)^{-1}N^{\mathsf T}AN=\operatorname{diag}(1,1,2,3),\\
\label{appA:eq:chi}
\chi_M(t)=(t-3)(t-2)(t-1)^2,\\
A_{21}=\begin{bmatrix}
0&0&0&0\\0&0&0&0\\-1&1&1&-1\\2&-2&-1&2\\
0&0&0&0\\0&0&0&0\\0&0&0&0\\0&0&0&0
\end{bmatrix}.
\end{gather}
Two rows of $A_{21}$ are nonzero; this is the map that distinguishes eigen from non-eigen directions.
\subsection{The six directions}
\label{appA:six}
The eigenspace of $\mu=1$ is $E_1=\operatorname{span}\{e_1,e_2\}\subseteq\mathbb R^4$ in $K_0$
coordinates, so $\mathbb P(E_1)\cong\mathbb P^1$; $(1,1)$ is the unique eigen direction
($A_{21}(c_1,c_2,0,0)^{\mathsf T}=0\iff c_1=c_2$), and its fibre $\mathbb P^3$ has one parameter more
than the fibres $\mathbb P^2$ of the other directions. The six directions and their verdicts are listed
in Table~\ref{appA:tab:dirs}.
\subsection{Chart variables, and the certificates in expanded form}
\label{appA:certificates}
For a non-eigen direction the fibre is $\mathbb P^2$ with homogeneous coordinates $(x_0\!:\!x_1\!:\!x_2)$;
for the eigen direction $(1\!:\!1)$ it is $\mathbb P^3$ with $(x_0\!:\!x_1\!:\!x_2\!:\!x_3)$. Below,
every maximal minor is written as a polynomial in $\lambda$ whose coefficients are polynomials in those
coordinates with integer coefficients, so that each identity can be checked by expansion. All
coefficients lie in $\mathbb Q$.

\emph{Bad directions.} In each of the five non-eigen directions the maximal minors share the factor
\begin{equation}\label{appA:eq:g}
g(x,\lambda)=x_2(\lambda-1)(\lambda-3)=x_2\lambda^2-4x_2\lambda+3x_2 .
\end{equation}
For $(1\!:\!0)$ the two nonzero maximal minors are, explicitly,
\begin{align}
m_1(x,\lambda)&=(\lambda-1)(\lambda-2)(\lambda-3)=\lambda^3-6\lambda^2+11\lambda-6,
\label{appA:eq:m1}\\
m_2(x,\lambda)&=x_2(\lambda-1)(\lambda-3)=x_2\lambda^2-4x_2\lambda+3x_2,
\label{appA:eq:m2}
\end{align}
and division in $\mathbb Q(x)[\lambda]$ gives, with zero remainders in both cases,
\begin{equation}\label{appA:eq:divcert}
m_1=g\cdot\frac{\lambda-2}{x_2},\qquad m_2=g\cdot 1 .
\end{equation}
For the remaining non-eigen directions the same factor $g$ occurs; the counts and the outcome of the
divisions are collected in Table~\ref{appA:tab:cert}.

\begin{table}[t]
\centering
\caption{Divisibility certificate per direction. ``Minors'' is the number of maximal minors that are not
identically zero; ``division'' records whether $g\mid m_I$ with zero remainder for every one of them.}
\label{appA:tab:cert}\label{appA:tab:dirs}
\begin{tabular}{@{}c c c c l@{}}
\toprule
Direction & Minors & $\deg_\lambda g$ & Division & Verdict \\
\midrule
$(1,0)$  & 2 & 2 & yes & entirely bad \\
$(0,1)$  & 2 & 2 & yes & entirely bad \\
$(1,1)$  & 8 & 0 & --- ($\gcd=1$) & good at $\xi_0$ \\
$(1,2)$  & 4 & 2 & yes & entirely bad \\
$(2,-1)$ & 4 & 2 & yes & entirely bad \\
$(3,1)$  & 4 & 2 & yes & entirely bad \\
\bottomrule
\end{tabular}
\end{table}

\emph{The good direction.} At the rational point
\begin{equation}\label{appA:eq:xipoint}
\xi_0=(1,2,3,4)\in\mathbb P^3
\end{equation}
of the fibre over $(1\!:\!1)$ the eight nonzero maximal minors specialize to
\begin{gather}
-(\lambda-1)(\lambda-2)(\lambda-3),\quad +(\lambda-1)(\lambda-2)(\lambda-3),\notag\\
5(\lambda-1)(\lambda-2),\quad -5(\lambda-1)(\lambda-2),\notag\\
(\lambda-1)(\lambda-3),\quad -(\lambda-1)(\lambda-3),\notag\\
-5(\lambda-2)(\lambda-3),\quad 5(\lambda-2)(\lambda-3),\label{appA:eq:spec}
\end{gather}
in the order $m_1,\dots,m_8$, and the Bezout identity
\begin{equation}\label{appA:eq:bezout}
\frac{\lambda-4}{10}\,m_4+\frac{\lambda-4}{2}\,m_5-\frac{1}{10}\,m_7=1
\end{equation}
holds. Its left-hand side is
$\frac{1}{2}\bigl[-(\lambda-4)(\lambda-1)(\lambda-2)+(\lambda-4)(\lambda-1)(\lambda-3)+(\lambda-2)(\lambda-3)\bigr]
=\frac{1}{2}\bigl[-(\lambda^2-5\lambda+4)+(\lambda^2-5\lambda+6)\bigr]=1$,
so \eqref{appA:eq:bezout} is an identity in $\mathbb Q[\lambda]$. Consequently
$\gcd_I m_I(\xi_0,\lambda)=1$, and $\xi_0$ is a feasible point of that fibre.

\emph{What these certificates establish.} Identities \eqref{appA:eq:divcert} and
\eqref{appA:eq:bezout} are finite rational identities: expanding either side and cancelling verifies
them in any exact arithmetic system, with no algebraic extension and no rounding. Together with the
full-rank leading-coefficient lemma of Section~\ref{sec:decision} they give the two conclusions used in
Section~\ref{sec:counterexample}: the fibres over the five non-eigen directions are entirely bad, and
the fibre over $(1\!:\!1)$ contains a feasible design.
\subsection{Reproduction}
\label{appA:reproduction}
Compute $N$, $M$, $\chi_M$, $A_{21}$ from $A$ and $C$ by rational elimination; for each direction build
$D_v$, $\mathcal Q_v$ and the maximal minors of $P(\lambda)=\lambda H-H_A$; divide them by $g$ in
$\mathbb Q(x)[\lambda]$ (fraction-free subresultants) to reproduce Table~\ref{appA:tab:cert}; substitute
$\xi_0$ and expand to verify \eqref{appA:eq:spec} and \eqref{appA:eq:bezout}.
